\documentclass[11pt,a4paper,oneside]{book}
\usepackage{mathscommun}
\setniveau{Seconde}
\renewcommand{\ReferenceProgramme}{Programme français de seconde GT -- rentrée 2026}
\renewcommand{\EditionDocument}{Édition du 18 octobre 2026}
\hypersetup{pdftitle={Mathématiques de seconde - Corrigés expliqués}}
\begin{document}
\titrecollection{Corrigés expliqués}{Comprendre la démarche et corriger ses erreurs.}{Cent quarante-quatre exercices, six par chapitre, et six problèmes de synthèse. Les approfondissements sont signalés. La numérotation correspond à celle du cours et de l'autre cahier.}
\frontmatter
\chapter*{Mode d'emploi}
Ces corrigés donnent une démarche complète, mais d'autres solutions correctes existent. Avant de lire, noter le point où la recherche s'est arrêtée. Après lecture, refermer le document et refaire la question ; une correction comprise doit pouvoir être reconstruite.

Les valeurs approchées sont indiquées comme telles. Les domaines, les valeurs interdites et les hypothèses des résultats font partie de la solution. Les exercices facultatifs gardent ce statut dans le corrigé.
\tableofcontents
\mainmatter
\chapter{Langage, ensembles et preuves}
\section*{Exercice 1.1 -- Ensembles}
$A=\{2,4,6,8,10\}$, $B=\{3,6,9\}$. L'intersection est $\{6\}$ (cardinal $1$), la réunion $\{2,3,4,6,8,9,10\}$ (cardinal $7$), le complémentaire $\{1,3,5,7,9\}$ (cardinal $5$). Le complémentaire est relatif à $E$.

\section*{Exercice 1.2 -- Couples}
On obtient $(0,a),(0,b),(0,c),(1,a),(1,b),(1,c)$, soit six couples. L'ordre distingue $(0,a)$ et $(a,0)$ ; le second n'appartient même pas à $A\times B$.

\section*{Exercice 1.3 -- Négations}
Les négations sont $x\leq2$ ; « $x<0$ ou $x\geq5$ » ; « $x\ne-1$ et $x\ne1$ ». Pour sortir d'une intersection, il suffit de manquer une condition ; pour sortir d'une réunion, il faut manquer les deux.

\section*{Exercice 1.4 -- Réciproque}
Si $n=6k$ avec $k\in\Z$, alors $n=3(2k)$, donc $3$ divise $n$. La réciproque est fausse : $3$ est multiple de $3$ sans être multiple de $6$.

\section*{Exercice 1.5 -- Une égalité d'ensembles}
Un élément du membre de gauche est dans $A$ et dans au moins l'un de $B,C$ ; il appartient donc à $A\cap B$ ou à $A\cap C$. Réciproquement, appartenir à l'un de ces deux ensembles impose l'appartenance à $A$ et à $B\cup C$. Les deux inclusions sont établies.

\section*{Exercice 1.6 -- L'ordre des mots}
La première est vraie : choisir $y=-x$, qui dépend de $x$. La seconde est fausse : $x=0$ imposerait $y=0$, alors que $x=1$ imposerait $y=-1$. Un même $y$ ne satisfait pas les deux exigences.

\chapter{Entiers, divisibilité et fractions}
\section*{Exercice 2.1 -- Diviseurs}
Les diviseurs de $36$ sont $1,2,3,4,6,9,12,18,36$. On a $84=6\times14=7\times12$ ; $84=8\times10+4$ n'est pas divisible par $8$.

\section*{Exercice 2.2 -- Dernier multiple}
$250=13\times19+3$ ; le multiple recherché est $247$. Le suivant est $260>250$, et les multiples de $13$ sont croissants avec leur quotient entier.

\section*{Exercice 2.3 -- Fraction}
$630=2\times3^2\times5\times7$ et $462=2\times3\times7\times11$. Le PGCD vaut $42$, d'où $630/462=15/11$. Les entiers $15$ et $11$ sont premiers entre eux.

\section*{Exercice 2.4 -- Parité}
Pour $2k+1$ et $2\ell+1$, la somme vaut $2(k+\ell+1)$ ; le produit vaut $4k\ell+2k+2\ell+1=2(2k\ell+k+\ell)+1$. Les quantités entre parenthèses sont entières.

\section*{Exercice 2.5 -- Lots}
Le nombre de lots divise les deux effectifs ; son maximum est $\PGCD(96,144)=48$. Chaque lot contient $96/48=2$ rouges et $144/48=3$ bleues.

\section*{Exercice 2.6 -- Un reste invariant}
Écrire $n=2k+1$ donne $n^2-1=4k(k+1)$. Parmi $k$ et $k+1$, l'un est pair, donc $k(k+1)=2m$ pour un entier $m$. Ainsi $n^2-1=8m$.

\chapter{Réels, intervalles et approximations}
\section*{Exercice 3.1 -- Classer}
Respectivement $\Z$, $\D$ car $7/8=0{,}875$, $\Q$ car $2/3$ n'a pas d'écriture décimale finie, $\N$ car $\sqrt{25}=5$, et $\R$ pour l'irrationnel $\sqrt3$. Pour ce dernier point, si $\sqrt3=p/q$ sous forme irréductible, alors $p^2=3q^2$. L'unicité de la décomposition en facteurs premiers impose que $3$ divise p ; écrire $p=3k$ donne $q^2=3k^2$, donc $3$ divise aussi q, contradiction.

\section*{Exercice 3.2 -- Intervalles}
$I\cap J=]0;2[$ ; $I\cup J=[-3;5]$ ; $\R\setminus I=]-\infty;-3[\cup[2;+\infty[$. La borne $2$ est exclue de $I$, donc incluse dans son complémentaire.

\section*{Exercice 3.3 -- Distance}
Le centre est $-2$ et le rayon $3$ : $x\in[-5;1]$. Pour le second intervalle, le centre est $7$ et le rayon $3$ : $|x-7|\leq3$.

\section*{Exercice 3.4 -- Encadrement}
$1{,}73^2=2{,}9929<3<3{,}0276=1{,}74^2$, donc $1{,}73<\sqrt3<1{,}74$. Pour l'arrondi, vérifier $1{,}735^2=3{,}010225>3$ : la racine est sous le milieu, donc l'arrondi est $1{,}73$. Le seul encadrement entre centièmes n'aurait pas suffi à choisir le plus proche.

\section*{Exercice 3.5 -- Irrationalité}
Si $1+\sqrt2$ était rationnel, lui soustraire le rationnel $1$ donnerait un rationnel. On obtiendrait que $\sqrt2$ est rationnel, contradiction. On utilise la stabilité de $\Q$ par soustraction.

\section*{Exercice 3.6 -- Chiffres significatifs}
On obtient $0{,}00457$ et $9{,}9\times10^4=99000$. Les pas d'arrondi sont respectivement $10^{-5}$ et $1000$ ; les erreurs absolues sont au plus $5\times10^{-6}$ et $500$. Ici les erreurs effectives valent $0{,}000003$ et $235$.

\chapter{Puissances et racines carrées}
\section*{Exercice 4.1 -- Puissances}
Les exposants donnent $4-7+1=-2$, donc $1/4$ ; puis $(-2)\times3+5=-1$, donc $1/3$. Les bases sont non nulles, ce qui autorise les exposants négatifs.

\section*{Exercice 4.2 -- Notation scientifique}
Le coefficient vaut $6\times4/3=8$ et l'exposant $-3+7-2=2$. Le résultat est $8\times10^2=800$.

\section*{Exercice 4.3 -- Racines}
La première expression vaut $5\sqrt2+4\sqrt2-3\sqrt2=6\sqrt2$. La seconde vaut $|-7|=7$, car une racine carrée est non négative.

\section*{Exercice 4.4 -- Comparer}
Les deux nombres du premier couple sont positifs ; leurs carrés valent $28$ et $25$, donc $2\sqrt7>5$. Pour le second, le carré de la somme vaut $9+2\sqrt{14}>9$, donc $\sqrt2+\sqrt7>3=\sqrt9$.

\section*{Exercice 4.5 -- Seuil}
$1{,}5^5=7{,}59375<10$ et $1{,}5^6=11{,}390625>10$. La suite des puissances est strictement croissante puisque le facteur est supérieur à $1$. Le premier rang est donc $6$.

\section*{Exercice 4.6 -- Une preuve}
Le quotient de racines est non négatif et son carré vaut $a/b$ ; l'unicité de la racine donne l'égalité. Pour $a=b=-1$, le quotient $a/b=1$ possède bien une racine, mais $\sqrt{-1}$ n'est pas un réel : le membre de droite n'existe pas dans $\R$.

\chapter{Expressions, identités et formules}
\section*{Exercice 5.1 -- Développer}
$(2x-3)^2=4x^2-12x+9$ ; $(x+4)(x-4)=x^2-16$. Dans le premier calcul, le terme croisé est $-2\times2x\times3$.

\section*{Exercice 5.2 -- Factoriser}
$6x^2-15x=3x(2x-5)$ ; la seconde expression vaut $(x+1)(2x-3+4)=(x+1)(2x+1)$. On met d'abord le facteur commun en évidence.

\section*{Exercice 5.3 -- Domaine conservé}
Le domaine est $\R\setminus\{-3\}$. Pour ces valeurs, la factorisation $(x-3)(x+3)$ donne $x-3$. La simplification ne supprime pas l'interdiction initiale : on ne peut pas évaluer le quotient en $-3$.

\section*{Exercice 5.4 -- Isoler}
$h=V/(\pi r^2)$ puisque le dénominateur est non nul. Puis $r^2=V/(\pi h)$ ; la condition de longueur sélectionne $r=\sqrt{V/(\pi h)}$. Sans $r\geq0$, il faudrait considérer les deux signes lorsque $V>0$.

\section*{Exercice 5.5 -- Choisir la forme}
$x^2-8x+20=(x-4)^2+4$. Son minimum vaut $4$, atteint en $x=4$. L'égalité à $8$ donne $(x-4)^2=4$, donc $x=2$ ou $x=6$.

\section*{Exercice 5.6 -- Moyennes}
$(a-b)^2\geq0$ donne $(a+b)^2\geq4ab$. Comme $a+b=12$, on obtient $144\geq4ab$, donc $ab\leq36$. L'égalité impose $a=b=6$.

\chapter{Équations, inéquations et signes}
\section*{Exercice 6.1 -- Premier degré}
La première donne $3x=15$, donc $S=\{5\}$. La seconde conduit à $-3=2$, impossible : $S=\varnothing$.

\section*{Exercice 6.2 -- Produit nul}
Le produit s'annule si $2x-1=0$ ou $x+4=0$, donc pour $x=1/2$ ou $x=-4$. Pour le carré, $x=\pm\sqrt{12}=\pm2\sqrt3$.

\section*{Exercice 6.3 -- Quotient}
Dans les deux cas, $x\ne1$. La première donne $x+3=2x-2$, donc $x=5$, admissible. La seconde se simplifie en $x+1=2$ sur le domaine, ce qui donne le seul candidat $1$, interdit. Son ensemble de solutions est vide.

\section*{Exercice 6.4 -- Signes}
Le produit est positif à l'extérieur des racines : $]-\infty,-2[\cup]3,+\infty[$. Le quotient est négatif entre $-2$ et $3$, nul en $3$ et interdit en $-2$ : solution $]-2,3]$.

\section*{Exercice 6.5 -- Budget}
Pour $n\in\N$, $14+3n\leq50$ équivaut à $n\leq12$. Le maximum est $12$ séances, pour un coût exactement égal à $50$ euros.

\section*{Exercice 6.6 -- Équivalence ?}
Il a supposé implicitement $x\ne0$ ; la division n'est donc pas équivalente sur tout le domaine. Le produit nul donne les deux solutions $0$ et $2$. Une autre méthode consiste à traiter d'abord $x=0$, puis le cas $x\ne0$.

\chapter{Repérage, distances et lieux géométriques}
\section*{Exercice 7.1 -- Distance et milieu}
$\overrightarrow{AB}$ a pour différences de coordonnées $(6,8)$, donc $AB=\sqrt{36+64}=10$. Le milieu est $I(2,6)$, obtenu en faisant la moyenne de chaque paire de coordonnées.

\section*{Exercice 7.2 -- Triangle rectangle}
$AB=4$, $AC=3$, $BC=\sqrt{16+9}=5$. Comme $AB^2+AC^2=BC^2$, la réciproque de Pythagore donne un angle droit en $A$. Cela se lit aussi sur les axes perpendiculaires.

\section*{Exercice 7.3 -- Cercle}
Le centre est $(2,-1)$ et le rayon $5$. Pour $(5,3)$, la somme vaut $3^2+4^2=25$, donc le point est sur le cercle.

\section*{Exercice 7.4 -- Médiatrice}
L'égalité des carrés des distances donne $(x+2)^2+y^2=(x-4)^2+y^2$, soit $12x=12$. Le lieu est $x=1$, droite verticale passant par le milieu de $[AB]$, donc sa médiatrice.

\section*{Exercice 7.5 -- Distance à un axe}
Les distances sont $4$ à l'axe horizontal et $3$ au vertical. Pour $P(t,0)$, $MP^2=(t-3)^2+16\geq16$, avec égalité uniquement si $t=3$.

\section*{Exercice 7.6 -- Rectangle général}
Prendre $(0,0),(a,0),(a,b),(0,b)$ dans un repère orthonormé. Les deux milieux valent $(a/2,b/2)$ ; chaque diagonale a pour carré de longueur $a^2+b^2$. On a traité tous les rectangles avec $a,b>0$.

\chapter{Vecteurs, translations et combinaisons}
\section*{Exercice 8.1 -- Déplacements}
$\overrightarrow{AB}=(3,4)$. On ajoute ce déplacement à $C$, donc $D(-1,6)$.

\section*{Exercice 8.2 -- Combinaison}
$3\vec u-2\vec v=(6,-9)-(-2,8)=(8,-17)$. Sa norme vaut $\sqrt{64+289}=\sqrt{353}$.

\section*{Exercice 8.3 -- Chasles}
Les deux premiers vecteurs donnent $\overrightarrow{AC}$. Puis $-\overrightarrow{DC}=\overrightarrow{CD}$, donc la somme vaut $\overrightarrow{AC}+\overrightarrow{CD}=\overrightarrow{AD}$.

\section*{Exercice 8.4 -- Parallélogramme}
L'égalité $\overrightarrow{AD}=\overrightarrow{BC}=(1,4)$ donne $D(2,5)$. Les diagonales ont alors le milieu commun $(3,3{,}5)$, ce qui fournit une vérification.

\section*{Exercice 8.5 -- Milieu vectoriel}
Par Chasles, la somme vaut $2\overrightarrow{MI}+\overrightarrow{IA}+\overrightarrow{IB}$. Les deux derniers vecteurs sont opposés puisque $I$ est le milieu. Il reste $2\overrightarrow{MI}$.

\section*{Exercice 8.6 -- Changement d'échelle}
La direction est conservée, le sens inversé et la norme vaut $(3/2)\times5=7{,}5$. Si $\vec u=(a,b)$, la nouvelle norme est $\sqrt{(9/4)(a^2+b^2)}=(3/2)\sqrt{a^2+b^2}$.

\chapter{Colinéarité, déterminant et bases}
\section*{Exercice 9.1 -- Déterminants}
Le premier vaut $2\times15-5\times6=0$ : colinéarité, facteur $3$. Le second vaut $14-15=-1\ne0$ : les vecteurs ne sont pas colinéaires.

\section*{Exercice 9.2 -- Alignement}
$\overrightarrow{AB}=(2,3)$ et $\overrightarrow{AC}=(4,6)=2\overrightarrow{AB}$. Comme $A\ne B$, cela prouve que $C$ appartient à $(AB)$.

\section*{Exercice 9.3 -- Paramètre}
Le déterminant vaut $3t-12$. Il est nul exactement pour $t=4$. Alors $(4,6)=2(2,3)$.

\section*{Exercice 9.4 -- Coordonnée nulle}
Le déterminant vaut $0\times(-7)-4\times0=0$. On peut aussi écrire $(0,-7)=(-7/4)(0,4)$ ; on divise par la coordonnée non nulle, pas par $0$.

\section*{Exercice 9.5 -- Dans une autre base}
Le système est $x+2y=7$, $2x-y=1$. De $y=2x-1$, on tire $5x=9$, donc $x=9/5$ et $y=13/5$. La base est valide car son déterminant est $-1-4=-5\ne0$.

\section*{Exercice 9.6 -- Parallèles ou confondues}
Les directeurs $(2,1)$ et $(4,2)$ sont colinéaires. Mais $C$ n'appartient pas à $(AB)$ : $\det((2,1),(0,3))=6\ne0$. Les droites sont parallèles distinctes.

\chapter{Droites et systèmes linéaires}
\section*{Exercice 10.1 -- Deux points}
Sa pente vaut $(9-3)/(4-1)=2$, puis $p=3-2=1$. Une équation est $y=2x+1$, ou $2x-y+1=0$.

\section*{Exercice 10.2 -- Point et directeur}
Le déterminant $3(y+1)-2(x-2)=0$ donne $-2x+3y+7=0$. Le point vérifie l'équation et un directeur est bien $(3,2)$.

\section*{Exercice 10.3 -- Verticale}
Elle a pour équation $x=4$, et pour directeur $(0,1)$. La différence des abscisses est nulle : le quotient de pente est indéfini. Ce n'est pas la courbe d'une fonction de $x$.

\section*{Exercice 10.4 -- Intersection}
L'égalité $3x-2=-2x+8$ donne $5x=10$, donc $x=2$, puis $y=4$. L'intersection est $(2,4)$.

\section*{Exercice 10.5 -- Trois systèmes}
Le premier système décrit la même droite deux fois : infinité de couples $(x,3-x)$. Le deuxième est contradictoire, car doubler la première ligne imposerait $6=7$. Dans le troisième, l'addition donne $2x=4$, donc l'unique couple est $(2,1)$.

\section*{Exercice 10.6 -- Région admissible}
C'est l'intersection du premier quadrant et du demi-plan sous $y=6-2x$, frontière comprise. Le triangle plein a pour sommets $(0,0),(3,0),(0,6)$. Les intersections avec les axes se trouvent en annulant successivement $y$ et $x$.

\chapter{Configurations et preuves géométriques}
\section*{Exercice 11.1 -- Milieux}
$\overrightarrow{MN}=\overrightarrow{AN}-\overrightarrow{AM}=\tfrac12(\overrightarrow{AC}-\overrightarrow{AB})=\tfrac12\overrightarrow{BC}$. On obtient le parallélisme et le rapport de longueurs $1/2$.

\section*{Exercice 11.2 -- Barycentre}
La somme des poids vaut $3$. Les coordonnées pondérées donnent $G((0+12)/3,(2-2)/3)=(4,0)$. Vérification : $\overrightarrow{GA}=(-4,2)$ et $2\overrightarrow{GB}=(4,-2)$ somment au vecteur nul.

\section*{Exercice 11.3 -- Médianes}
L'isobarycentre est $G(2,3)$. Le milieu de $[BC]$ est $I(3,4{,}5)$ ; $\overrightarrow{AG}=(2/3)\overrightarrow{AI}$, donc $G$ est aux deux tiers depuis $A$.

\section*{Exercice 11.4 -- Cercle circonscrit}
Le triangle est rectangle en $A$. Les médiatrices des côtés sur les axes sont $x=2$ et $y=1{,}5$, donc le centre est $(2,1{,}5)$. Sa distance à $A$ vaut $\sqrt{4+2{,}25}=2{,}5$, moitié de l'hypoténuse $5$.

\section*{Exercice 11.5 -- Aire}
La hauteur associée au côté $5$ vaut $8\sin30\degre=4$. L'aire est donc $(5\times4)/2=10$. C'est $\frac12 ab\sin\theta$, obtenue à partir de base fois hauteur divisée par deux.

\section*{Exercice 11.6 -- Lieu}
L'égalité de distances équivaut à $x^2+(y-2)^2=y^2$, donc $y=x^2/4+1$. Pour $(2,2)$, la distance à $F$ vaut $2$, comme celle à l'axe ; il appartient bien à la parabole.

\chapter{Fonctions : définir, représenter, modéliser}
\section*{Exercice 12.1 -- Images et antécédents}
$f(-2)=9$. Les antécédents de $4$ vérifient $x-1=\pm2$, donc sont $-1$ et $3$. $0$ possède le seul antécédent $1$ ; $-1$ n'en possède aucun puisque la fonction est non négative.

\section*{Exercice 12.2 -- Domaine}
La racine exige $x\leq3$, le dénominateur $x\ne-2$. Le domaine est $]-\infty,-2[\cup]-2,3]$.

\section*{Exercice 12.3 -- Appartenance}
Pour $A$, $2^2+1=5$, donc oui. Pour $B$, $(-2)^2+1=5\ne-3$, donc non. Il faut remplacer l'abscisse dans la formule et comparer à l'ordonnée.

\section*{Exercice 12.4 -- Tableau insuffisant}
Non : $f(x)=0$ et $g(x)=x(x-1)(x+1)$ coïncident sur ces trois entrées, mais $g(2)=6\ne0=f(2)$. Un tableau fini ne détermine pas une fonction quelconque sur $\R$.

\section*{Exercice 12.5 -- Modèle géométrique}
Les côtés vérifient $2x+2y=30$, donc $y=15-x$. L'aire est $A(x)=x(15-x)$ pour $0<x<15$. Le domaine du modèle est plus restreint que celui du polynôme.

\section*{Exercice 12.6 -- Modèle discret}
$u(4)=18$. L'équation $8+2{,}5n=19$ donne $n=4{,}4$, qui n'est pas entier. Ce prix n'est donc pas réalisable dans le modèle, même s'il existe un point d'abscisse $4{,}4$ sur la droite prolongée.

\chapter{Fonctions affines et comparaisons}
\section*{Exercice 13.1 -- Reconstruire}
La pente vaut $(14-5)/(4-1)=3$. Puis $p=5-3\times1=2$, donc $f(x)=3x+2$. Les deux données sont vérifiées.

\section*{Exercice 13.2 -- Signe et variation}
La pente $-4$ est négative, donc $f$ est strictement décroissante sur $\R$. Son zéro est $3$ ; elle est positive pour $x<3$, nulle en $3$, négative pour $x>3$. Les deux études ne doivent pas être confondues.

\section*{Exercice 13.3 -- Comparer des tarifs}
$A-B=20-4x$. Les tarifs sont égaux pour $x=5$ ; $B$ est moins cher pour $0\leq x<5$, $A$ pour $x>5$. À $5$, chacun coûte $40$.

\section*{Exercice 13.4 -- Une preuve}
Pour $a<b$, $f(b)-f(a)=m(b-a)$. Les deux facteurs étant positifs, la différence l'est aussi, donc $f(a)<f(b)$. Les entrées ont été choisies quelconques : la preuve porte sur tout le domaine.

\section*{Exercice 13.5 -- Interpolation}
La pente est $(30-18)/(8-2)=2$, d'où $f(t)=18+2(t-2)$. À $t=5$, elle vaut $24$. Le résultat suppose une variation affine entre les relevés ; les deux mesures seules ne l'imposent pas.

\section*{Exercice 13.6 -- Modèle physique}
$V(t)=75-2{,}5t$ sur $[0;30]$. L'équation $75-2{,}5t=20$ donne $t=22$ minutes, valeur admissible. Le modèle s'arrête à la vidange, $t=30$.

\chapter{Valeur absolue et fonction carré}
\section*{Exercice 14.1 -- Distances}
L'égalité donne $x=-3$ ou $x=5$. L'inégalité décrit tous les points à distance au plus $4$ du centre $1$, donc $[-3,5]$.

\section*{Exercice 14.2 -- Carrés}
La première donne $]-4,4[$. La seconde donne $]-\infty,-\sqrt5]\cup[\sqrt5,+\infty[$. Les bornes sont incluses seulement pour l'inégalité large.

\section*{Exercice 14.3 -- Comparer}
La différence $x^2-x=x(x-1)$ est négative pour $0<x<1$, donc $0{,}7^2<0{,}7$. Elle est positive pour $x>1$, donc $1{,}3^2>1{,}3$.

\section*{Exercice 14.4 -- Parabole déplacée}
Le carré est minimal lorsque $x+2=0$, donc en $x=-2$, avec valeur $-3$. La fonction décroît sur $]-\infty,-2]$ et croît sur $[-2,+\infty[$, par translation de la fonction carré.

\section*{Exercice 14.5 -- Solutions et domaine}
On obtient $(x+2)^2=9$, donc $x+2=-3$ ou $3$ : les solutions réelles sont $-5$ et $1$. Sur $[0,+\infty[$, seule $1$ est admissible.

\section*{Exercice 14.6 -- Un contre-exemple}
$-3<-2$, mais $9>4$. Avec $0\leq a<b$, on a $b^2-a^2=(b-a)(b+a)>0$. La restriction aux positifs rend l'affirmation correcte.

\chapter{Inverse, racine carrée et cube}
\section*{Exercice 15.1 -- Inverse}
La première donne $x=1/4$. Pour la seconde, étudier $(1+2x)/x\leq0$ : le numérateur s'annule en $-1/2$, le dénominateur en $0$ interdit. Le quotient est négatif ou nul sur $[-1/2,0[$.

\section*{Exercice 15.2 -- Racine}
Le domaine exige $x\geq1$. L'égalité donne $x-1=16$, soit $17$. L'inégalité donne $0\leq x-1<16$, donc $[1,17[$.

\section*{Exercice 15.3 -- Cube}
Le cube est strictement croissant sur $\R$. L'égalité donne $x=-3$ ; l'inégalité équivaut à $x<2$, soit $]-\infty,2[$.

\section*{Exercice 15.4 -- Comparer des images}
Sur $]-\infty,0[$, l'inverse décroît ; puisque $-4<-2$, $1/(-4)>1/(-2)$. La racine croît sur les positifs, donc $\sqrt{11}<\sqrt{12}$.

\section*{Exercice 15.5 -- Solution parasite}
Il faut $x+6\geq0$ et surtout $x\geq0$, puisque la racine est non négative. Sur $[0,+\infty[$, l'égalité équivaut à $x+6=x^2$, donc $(x-3)(x+2)=0$. Seul $3$ convient ; $-2$ est une solution de l'équation carrée, pas de l'initiale.

\section*{Exercice 15.6 -- Symétrie}
Un point du carré a la forme $(t,t^2)$ avec $t\geq0$. Son image est $(t^2,t)$ ; comme $\sqrt{t^2}=t$, il appartient à la courbe de la racine. Réciproquement, tout point $(x,\sqrt x)$ s'obtient avec $t=\sqrt x$.

\chapter{Variations, extrémums et optimisation}
\section*{Exercice 16.1 -- Lire des variations}
Le minimum vaut $-1$, atteint en $1$ ; le maximum vaut $8$, atteint en $-2$. Les deux autres entrées appartiennent à des intervalles de variation différents ; ces seules informations ne permettent pas de comparer leurs images.

\section*{Exercice 16.2 -- Borne non atteinte}
Elle vérifie $0<f(x)<4$. Les bornes ne sont pas atteintes et la fonction est strictement croissante. Elle n'a donc ni minimum ni maximum sur cet intervalle ouvert.

\section*{Exercice 16.3 -- Optimiser}
Ses côtés sont $x$ et $12-x$, avec $0<x<12$. L'aire vaut $x(12-x)=36-(x-6)^2\leq36$. L'égalité est atteinte pour $x=6$, admissible. Le carré de côté $6$ réalise l'aire maximale $36\unite{m^2}$.

\section*{Exercice 16.4 -- Dichotomie}
Après $n$ étapes, l'amplitude vaut $2^{-n}$. Or $2^6=64<100$ et $2^7=128\geq100$. Il faut au moins $7$ étapes ; l'amplitude vaut alors $1/128\simeq0{,}0078125$.

\section*{Exercice 16.5 -- Grille et preuve}
Les valeurs sont $0{,}91$, $0{,}96$, $0{,}51$ : la grille retient $0{,}96$ en $0{,}5$. Mais $f(x)\leq1$, avec égalité en $0{,}3$, donc le maximum global vaut $1$. Une grille peut manquer le point optimal.

\section*{Exercice 16.6 -- Longueur approchée}
Avec un segment, la longueur vaut $\sqrt{1^2+1^2}=\sqrt2$. Avec deux segments, les différences de coordonnées sont $(1/2,1/4)$ puis $(1/2,3/4)$, d'où $L_2=(\sqrt5+\sqrt{13})/4\simeq1{,}4604$. Ce sont des longueurs polygonales, pas une formule exacte de la longueur de la parabole.

\chapter{Proportions et évolutions}
\section*{Exercice 17.1 -- Une proportion}
$N=54/0{,}18=300$. Le calcul inverse utilise la proportion comme coefficient, et non comme effectif.

\section*{Exercice 17.2 -- Proportion de proportion}
$0{,}40\times0{,}30=0{,}12$, soit $12\,\%$. Le deuxième pourcentage porte sur le sous-groupe des usagers du bus.

\section*{Exercice 17.3 -- Successions}
Le coefficient global est $1{,}10\times0{,}85=0{,}935$. Le taux global est $-0{,}065$, soit une baisse de $6{,}5\,\%$.

\section*{Exercice 17.4 -- Réciproque}
Le coefficient initial est $0{,}7$. Il faut multiplier ensuite par $1/0{,}7=10/7$, soit une hausse de $3/7\simeq42{,}86\,\%$. Une hausse de $30\,\%$ ne suffirait pas.

\section*{Exercice 17.5 -- Points ou pourcentage}
Elle augmente de $3$ points de pourcentage. La hausse relative est $(15-12)/12=1/4=25\,\%$. Les deux réponses ne désignent pas la même grandeur.

\section*{Exercice 17.6 -- Pondérer}
Les effectifs oui sont $10$ et $30$, soit $40$ sur $100$ : $40\,\%$. La moyenne simple des deux pourcentages donnerait $37{,}5\,\%$, car elle attribuerait à tort le même poids aux deux groupes.

\chapter{Statistiques : position et lecture critique}
\section*{Exercice 18.1 -- Indicateurs}
La somme vaut $52$, donc la moyenne $5{,}2$. La médiane est $(5+5)/2=5$. Les rangs des quartiles sont $\lceil10/4\rceil=3$ et $\lceil30/4\rceil=8$, donc $Q_1=2$, $Q_3=8$. L'étendue vaut $9$, l'écart interquartile $6$.

\section*{Exercice 18.2 -- Moyenne de groupes}
Les sommes sont $12\times14=168$ et $18\times11=198$. Leur total est $366$ pour $30$ élèves, donc la moyenne vaut $12{,}2$.

\section*{Exercice 18.3 -- Transformation}
La linéarité donne $\overline y=3\times7-2=19$. On peut le retrouver en sommant $3x_i-2$ : la somme devient trois fois l'ancienne somme moins deux fois l'effectif.

\section*{Exercice 18.4 -- Ajout extrême}
La somme devient $152$, l'effectif $11$, donc la moyenne vaut $152/11\simeq13{,}82$. La valeur centrale est de rang $6$, donc la médiane reste $5$. La valeur extrême déplace beaucoup plus la moyenne.

\section*{Exercice 18.5 -- Lire un diagramme}
La fréquence vaut $45/180=1/4$. Les angles sont $360/4=90\degre$ et $180/4=45\degre$. La relation vient de la proportionnalité de l'angle à l'aire du secteur pour un rayon fixé.

\section*{Exercice 18.6 -- Mêmes centres ?}
La moyenne et la médiane valent $5$ dans les deux séries, mais les étendues sont $2$ et $10$. Les distributions ne sont pas identiques : $B$ est plus dispersée. Deux indicateurs centraux ne déterminent pas une série.

\chapter{Dispersion et séries regroupées en classes}
\section*{Exercice 19.1 -- Calculer une dispersion}
$\bar x=4$. Par les écarts, $V=(4+0+0+4)/4=2$. Par la moyenne des carrés, $V=(4+16+16+36)/4-4^2=18-16=2$. Ainsi $\sigma=\sqrt2\simeq1{,}414$. L'unité de $\sigma$ est celle des données ; celle de $V$ est son carré.

\section*{Exercice 19.2 -- Changer d'unité}
La moyenne devient $1{,}8\times18+32=64{,}4$. Chaque écart à la moyenne est multiplié par $1{,}8$ : $(1{,}8C_i+32)-(1{,}8\bar C+32)=1{,}8(C_i-\bar C)$. Donc l'écart-type est $1{,}8\times3=5{,}4$. Une translation ne change pas la dispersion.

\section*{Exercice 19.3 -- Classes et histogramme}
L'effectif total est $20$. Les hauteurs peuvent être les effectifs $4,8,8$, puisque les largeurs sont égales. Le polygone joint $(0;0)$, $(10;0{,}2)$, $(20;0{,}6)$ et $(30;1)$. Les centres donnent $\bar x\simeq(4\times5+8\times15+8\times25)/20=17$. La médiane estimée est dans la deuxième classe : $10+10(0{,}5-0{,}2)/(0{,}6-0{,}2)=17{,}5$. Ces nombres sont des estimations : les effectifs par classe ne déterminent pas les valeurs individuelles.

\section*{Exercice 19.4 -- Ce que le regroupement perd}
$\bar x=(4\times4+8\times14+8\times26)/20=336/20=16{,}8$. L'écart-type reste inconnu : dans la première classe, quatre $4$ et les valeurs $2,4,4,6$ donnent la même moyenne mais des dispersions différentes. On peut garder identiques toutes les autres observations. Les moyennes de classe ne suffisent donc pas à déterminer la moyenne des carrés.

\section*{Exercice 19.5 -- Choisir les indicateurs}
Dans les deux cas la médiane vaut $12$, $Q_1$ est la valeur de rang $2$, soit $11$, et $Q_3$ celle de rang $4$, soit $13$. L'écart interquartile vaut $2$. Les moyennes sont $12$ et $146/5=29{,}2$ ; les étendues sont $4$ et $90$. La valeur $100$ modifie fortement moyenne et étendue sans modifier ici médiane et quartiles. Il faut décrire à la fois le centre, la dispersion et les valeurs extrêmes.

\section*{Exercice 19.6 -- Décomposer une variance (facultatif)}
La moyenne globale est $5$. Pour chaque observation, écrire $x-5=(x-\bar x_g)+(\bar x_g-5)$ et développer le carré. Dans chaque groupe, la somme des $x-\bar x_g$ est nulle : le terme croisé disparaît. La variance totale est donc $(1+4)/2+((2-5)^2+(8-5)^2)/2=2{,}5+9=11{,}5$. La dispersion entre les deux moyennes doit s'ajouter à la dispersion interne.

\chapter{Tableaux croisés et fréquences conditionnelles}
\section*{Exercice 20.1 -- Compléter des effectifs}
\begin{center}\begin{tabular}{lrrr}\toprule & Casque & Sans casque & Total\\\midrule A & 90&30&120\\ B&40&40&80\\ Total&130&70&200\\\bottomrule\end{tabular}\end{center} Chaque marge est une somme ; le total général se retrouve dans les deux directions.

\section*{Exercice 20.2 -- Choisir le dénominateur}
Dans toute la population : $130/200=65\,\%$. Parmi A : $90/120=75\,\%$. Parmi les porteurs : $90/130=9/13\simeq69{,}23\,\%$. Les deux dernières fractions ont le même numérateur, mais décrivent des populations de référence différentes.

\section*{Exercice 20.3 -- Nominal ou ordinal}
Le canton est qualitatif nominal ; la satisfaction est qualitative ordinale ; le nombre d'enfants est quantitatif discret ; la taille est quantitative continue (mesurée à une précision donnée). Une moyenne est directement pertinente pour les deux derniers. Coder faible, moyenne, forte par $1,2,3$ ne prouve pas que les écarts de satisfaction sont égaux.

\section*{Exercice 20.4 -- Lire des données réelles}
Le canton compte $195\,091$ personnes. La première proportion est $97\,721/195\,091\simeq50{,}09\,\%$. La seconde est $97\,721/339\,161\simeq28{,}81\,\%$. Il s'agit de deux fréquences conditionnelles distinctes. La première décrit la composition du canton ; la seconde la répartition géographique de la population féminine. Ce sont des données datées, pas une estimation de la population présente aujourd'hui.

\section*{Exercice 20.5 -- Filtrer avec ET, OU, NON}
Intersection : $90$. Union : $120+130-90=160$, car les $90$ personnes de l'intersection seraient comptées deux fois. Complément de l'union : $200-160=40$, ce sont les personnes de B sans casque. Les mots ET, OU et NON correspondent à intersection, union et complémentaire.

\section*{Exercice 20.6 -- Une inversion trompeuse}
Non : la proportion de filles qui font du théâtre est $9/18=1/2$. La fraction $9/12=3/4$ signifie que les trois quarts des élèves qui font du théâtre sont des filles. Le tableau complet comporte $9$ filles sans théâtre, $3$ garçons avec théâtre et $9$ garçons sans théâtre.

\chapter{Probabilités finies et dénombrement}
\section*{Exercice 21.1 -- Construire une loi}
Écrire $P(k)=ck$. La somme vaut $c(1+2+3+4)=10c=1$, donc $c=1/10$. La loi est $(0{,}1;0{,}2;0{,}3;0{,}4)$ et $P(\text{pair})=0{,}2+0{,}4=0{,}6$. L'équiprobabilité ne s'applique pas ici.

\section*{Exercice 21.2 -- Deux dés}
Les $36$ couples ordonnés sont équiprobables sous les hypothèses annoncées. Il y a six couples de somme $7$, donc $1/6$. Pour au moins un $6$, le contraire a $5\times5=25$ couples : $P=1-25/36=11/36$. Les sommes ne sont pas équiprobables.

\section*{Exercice 21.3 -- Réunion et intersection}
Les régions $A\setminus B$, $A\cap B$ et $B\setminus A$ ont des probabilités $0{,}2$, $0{,}2$ et $0{,}3$. Leur somme est $0{,}7$. En général, additionner $P(A)$ et $P(B)$ compte deux fois l'intersection ; en retirer une copie donne $P(A\cup B)=P(A)+P(B)-P(A\cap B)$. Le contraire vaut $0{,}3$.

\section*{Exercice 21.4 -- Remettre ou ne pas remettre}
Avec remise et mélange à chaque fois, la probabilité est $(3/5)^2=9/25$. Sans remise, après une rouge il reste deux rouges parmi quatre boules : $(3/5)(2/4)=3/10$. On peut aussi compter $3\times2=6$ couples favorables parmi $5\times4=20$ couples ordonnés sans remise.

\section*{Exercice 21.5 -- Un modèle est une hypothèse}
$0{,}63$ est une fréquence observée. Elle n'est pas automatiquement la probabilité inconnue. Une simulation sous l'hypothèse d'une pièce équilibrée permet d'observer la variabilité des fréquences et de comparer l'observation au modèle. Elle ne prouve ni l'équilibrage ni une valeur exacte du paramètre ; l'étude formelle d'un test dépasse ce chapitre.

\section*{Exercice 21.6 -- Compter un complément}
Il y a $10^3=1000$ codes. Sans répétition : $10\times9\times8=720$. Donc $P(\text{au moins une répétition})=1-720/1000=0{,}28$. Le complément évite de compter plusieurs fois un code comme $777$.

\chapter{Probabilités conditionnelles et arbres}
\section*{Exercice 22.1 -- Du tableau à la probabilité}
$P(A)=120/200=0{,}6$ ; $P_A(C)=90/120=0{,}75$ ; $P_C(A)=90/130=9/13$ ; $P(A\cap C)=90/200=0{,}45$. On vérifie $P(A)P_A(C)=0{,}6\times0{,}75=0{,}45$.

\section*{Exercice 22.2 -- Deux chemins}
Première branche : A $0{,}6$, B $0{,}4$. Deuxième niveau : défaut/non-défaut, respectivement $(0{,}02;0{,}98)$ dans A et $(0{,}05;0{,}95)$ dans B. Les deux chemins disjoints vers un défaut ont des probabilités $0{,}012$ et $0{,}020$. Leur somme est $0{,}032$. Parmi les défauts, la part venant de A vaut $0{,}012/0{,}032=3/8=37{,}5\,\%$.

\section*{Exercice 22.3 -- Un diagnostic fictif}
\begin{center}\begin{tabular}{lrrr}\toprule & Positif & Négatif & Total\\\midrule Malade&90&10&100\\ Non-malade&495&9405&9900\\ Total&585&9415&10000\\\bottomrule\end{tabular}\end{center}
Non : parmi les positifs, la proportion de malades vaut $90/585=2/13\simeq15{,}38\,\%$. $90\,\%$ est la probabilité d'un positif \emph{sachant que} la personne est malade. Ces paramètres sont inventés pour l'exercice et ne décrivent aucun test réel.

\section*{Exercice 22.4 -- Sensibilité et spécificité}
Sensibilité $P_M(T)=90\,\%$ ; spécificité $P_{\overline M}(\overline T)=95\,\%$ ; taux de faux négatifs parmi les malades $10\,\%$ ; taux de faux positifs parmi les non-malades $5\,\%$. Dans toute la population, $P(\overline M\cap T)=0{,}99\times0{,}05=0{,}0495$, soit $4{,}95\,\%$. Le dénominateur doit toujours être précisé.

\section*{Exercice 22.5 -- Indépendance (facultatif)}
$P(A)P(B)=0{,}2=P(A\cap B)$ : ils sont indépendants. On retrouve $P_A(B)=0{,}2/0{,}4=0{,}5=P(B)$. Ils ne sont pas incompatibles puisque leur intersection a une probabilité positive. Indépendance ne signifie pas impossibilité de se produire ensemble.

\section*{Exercice 22.6 -- Pourquoi la multiplication ?}
La définition est $P_A(B)=P(A\cap B)/P(A)$. Multiplier par $P(A)$ donne $P(A\cap B)=P(A)P_A(B)$. Si $P(A)=0$, le quotient n'est pas défini dans ce cadre élémentaire. L'intersection a toutefois une probabilité nulle puisqu'elle est incluse dans A. Il ne faut pas écrire une probabilité conditionnelle qui n'existe pas.

\chapter{Algorithmique et Python}
\section*{Exercice 23.1 -- Tracer des variables}
Les valeurs de $k$ sont $1,2,3,4$, et celles de \texttt{s} après addition sont $1,5,14,30$. La borne finale de \texttt{range} est exclue. Avec \texttt{range(1, 4)}, on s'arrête à $k=3$ et le résultat vaut $14$.

\section*{Exercice 23.2 -- Dernier multiple}
\begin{lstlisting}[language=Python]
def dernier_multiple(a, b):
    if a <= 0 or b < 0:
        raise ValueError("a > 0 et b >= 0")
    return a * (b // a)
\end{lstlisting}
La division euclidienne donne $b=aq+r$ avec $0\le r<a$ ; $aq\le b<a(q+1)$. Pour $a=7$, $b=21$, on obtient $21$, et pour $b=20$, on obtient $14$. Le test sur un multiple exact vérifie le signe « inférieur ou égal ». Les entrées sont supposées entières.

\section*{Exercice 23.3 -- Un encadrement certifié}
$1414^2=1\,999\,396<2\,000\,000$ et $1415^2=2\,002\,225>2\,000\,000$. Donc $k=1414$ et $1{,}414<\sqrt2<1{,}415$. Les produits d'entiers Python sont exacts ; on évite de faire dépendre la décision finale de l'arrondi d'un flottant. La largeur de l'encadrement vaut $10^{-3}$.

\section*{Exercice 23.4 -- Lire une fonction statistique}
La liste doit être non vide et contenir des nombres. \texttt{n} est l'effectif, \texttt{m} la moyenne. On additionne les carrés des écarts, on divise par $n$, puis on prend la racine : le résultat est l'écart-type de population. La variance est l'expression avant la racine. L'écriture avec une liste est un support de lecture ; la maîtrise générale des listes n'est pas exigée en seconde.

\section*{Exercice 23.5 -- Simuler sans conclure trop vite}
\begin{lstlisting}[language=Python]
from random import random
def frequence_face(n):
    if n <= 0:
        raise ValueError("n doit etre positif")
    succes = 0
    for _ in range(n):
        if random() < 0.5:
            succes += 1
    return succes / n
\end{lstlisting}
Les appels utilisent des tirages pseudo-aléatoires différents. Les fréquences fluctuent même avec le même nombre de lancers. Le code simule un modèle d'indépendance et de probabilité $1/2$ ; une sortie particulière ne prouve pas une loi générale.

\section*{Exercice 23.6 -- Longueur polygonale}
Poser $x_k=k/n$ et additionner, pour $k=0,\ldots,n-1$, les distances entre $(x_k;x_k^2)$ et $(x_{k+1};x_{k+1}^2)$. Chaque distance vaut $\sqrt{(1/n)^2+(x_{k+1}^2-x_k^2)^2}$. Pour $n=1$, on obtient $\sqrt2$. Un tableau numérique ne démontre ni la convergence, ni une erreur maximale ; ces questions relèvent d'une étude d'analyse supplémentaire.

\chapter{Approfondir les preuves et préparer la suite}
\section*{Exercice 24.1 -- L'ordre des quantificateurs}
La première est vraie : choisir $y=x+1$. La seconde est fausse : pour un $y$ proposé, prendre $x=y$ (ou $y+1$). Les négations sont respectivement : « il existe un réel $x$ tel que, pour tout réel $y$, $y\le x$ » ; « pour tout réel $y$, il existe un réel $x$ tel que $y\le x$ ». La dépendance de $y$ à $x$ est décisive.

\section*{Exercice 24.2 -- Une somme télescopique}
$k^2-(k-1)^2=2k-1$. En additionnant pour $k=1,\ldots,n$, les termes intermédiaires s'annulent : la somme vaut $n^2-0^2=n^2$. Pour $n=4$, $1+3+5+7=16$. Ce calcul particulier illustre la formule ; seule l'identité sommée traite tout entier $n\ge1$.

\section*{Exercice 24.3 -- Une récurrence guidée}
Au rang $0$, les deux membres valent $1$. Si $(1+u)^n\ge1+nu$, multiplier par $1+u>0$ donne $(1+u)^{n+1}\ge1+(n+1)u+nu^2\ge1+(n+1)u$. L'hérédité est établie ; le principe de récurrence donne la propriété à tout rang. Choisir un entier $n>(M-1)/u$ rend $1+nu>M$. L'existence d'un entier assez grand utilise la propriété archimédienne des réels.

\section*{Exercice 24.4 -- Une borne de fluctuation}
$4n\varepsilon^2=4\times10\,000\times0{,}0025=100$ ; la borne vaut $0{,}01$. La probabilité d'un écart d'au moins $0{,}05$ est donc au plus $1\,\%$ sous les hypothèses du modèle. Ce n'est pas une impossibilité. La borne peut être peu précise ; elle justifie un rapprochement en probabilité, pas une garantie sur chaque tirage.

\section*{Exercice 24.5 -- Trouver le trou d'une preuve}
Le passage de $(a-b)(a+b)=b(a-b)$ à $a+b=b$ divise par $a-b$. Or $a=b$, donc $a-b=0$ : cette division est interdite. Les égalités précédentes sont correctes mais deviennent $0=0$, qui ne permet aucune simplification par le facteur nul.

\section*{Exercice 24.6 -- Démontrer seul}
$(\sqrt a-\sqrt b)^2\ge0$ donne $a+b\ge2\sqrt{ab}$, avec égalité si et seulement si $a=b$. En élevant au carré des quantités positives, $(a+b)^2\ge4ab$. En divisant par $ab(a+b)>0$, on obtient $(a+b)/(ab)\ge4/(a+b)$, soit l'inégalité demandée. Le cas d'égalité reste $a=b$. Le contrôle de la positivité justifie chaque opération.

\chapter{Six problèmes de synthèse}
\section*{Problème 1 -- Un jardin et son coût}
La longueur vaut $40-2x$ ; un rectangle non dégénéré exige $0<x<20$. L'aire vaut $x(40-2x)=-2x^2+40x=200-2(x-10)^2$. Elle est au plus $200$, atteinte pour $x=10$ et une longueur $20$.
L'inégalité $A(x)\ge150$ équivaut à $(x-10)^2\le25$, donc $x\in[5;15]$, inclus dans le domaine géométrique.
Pour $40$ mètres, les devis valent $720$ et $680$ euros. Leur différence est $-80+3L$ : le premier est moins cher pour $0\le L<80/3$, ils sont égaux à $80/3$, et le second est moins cher au-delà. Toutes les aires sont exprimées en mètres carrés et les coûts en euros.

\section*{Problème 2 -- Un triangle sous plusieurs regards}
$AB=6$, $AC=\sqrt{20}=2\sqrt5$, $BC=\sqrt{32}=4\sqrt2$. La hauteur issue de C vaut $4$, donc l'aire est $12$.
$I(4;2)$ ; AI a pour équation $x-2y=0$. G la vérifie, et $\overrightarrow{AG}=(2/3)\overrightarrow{AI}$.
Le milieu de AC est $J(1;2)$. $\overrightarrow{BG}=(-10/3;4/3)=(2/3)(-5;2)=(2/3)\overrightarrow{BJ}$. G est aux deux tiers de chaque médiane en partant du sommet.
La médiatrice de AB est $x=3$. L'égalité $OA^2=OC^2$ donne $x^2+y^2=(x-2)^2+(y-4)^2$, soit $x+2y=5$. Ainsi $O(3;1)$ et $R=\sqrt{10}$. On vérifie $OB^2=OC^2=10$.
La hauteur issue de C est $x=2$. BC a une pente $-1$ ; sa perpendiculaire par A a pour équation $y=x$. Elles se coupent en $H(2;2)$, qui appartient aussi à la hauteur issue de B, de pente $-1/2$.

\section*{Problème 3 -- Un tarif et une expérience}
$m(n)=80/n+0{,}25$. L'inverse est strictement décroissant sur les positifs, donc le coût moyen décroît avec $n$.
$80/n+0{,}25\le1$ équivaut, puisque $n>0$, à $80\le0{,}75n$, soit $n\ge320/3\simeq106{,}667$. Le plus petit entier est $107$. La courbe continue sert d'aide à la lecture ; seules ses abscisses entières positives représentent un nombre d'utilisations.
\begin{lstlisting}[language=Python]
n = 1
while 80 / n + 0.25 > 1:
    n += 1
\end{lstlisting}
L'entier $107$ vérifie le seuil, et la boucle visite successivement tous les entiers de $1$ à cette valeur : elle termine et son premier arrêt est minimal. Pour éviter un arrondi près de la frontière, on peut tester exactement \texttt{while 320 > 3 * n}.

\section*{Problème 4 -- Comparer deux distributions}
Dans A, moyenne et médiane valent $10$, $Q_1=8$, $Q_3=12$ et $\sigma=\sqrt{(10\times4+10\times4)/20}=2$. La boîte va de $8$ à $12$, la médiane est $10$, et les moustaches sont aux extrêmes $8$ et $12$.
Dans B, moyenne et médiane valent aussi $10$. $Q_1$ est de rang $5$, soit $4$ ; $Q_3$ de rang $15$, soit $10$. L'écart-type vaut $\sqrt{(5\times36+10\times0+5\times36)/20}=\sqrt{18}=3\sqrt2$. La boîte est $[4;10]$ avec médiane $10$, et maximum $16$. B est plus dispersé selon l'écart-type et l'écart interquartile ($6$ contre $4$). La convention des quartiles explique l'asymétrie de sa boîte.
Après modification, A a moyenne $12$ et écart-type $2$ ; B a moyenne $11$ et écart-type $3{,}3\sqrt2$. Un ajout fixe donne le même gain absolu ; une hausse en pourcentage donne ici $0{,}4$, $1$ ou $1{,}6$ point selon la note initiale. Les deux transformations n'ont pas le même effet.

\section*{Problème 5 -- Du contrôle qualité à la décision}
Les chemins donnent : défectueuse rejetée $0{,}03\times0{,}95=0{,}0285$ ; défectueuse acceptée $0{,}0015$ ; conforme rejetée $0{,}97\times0{,}04=0{,}0388$ ; conforme acceptée $0{,}9312$. Le tableau théorique comporte donc $2850$, $150$, $3880$ et $93\,120$ pièces.
La probabilité de rejet vaut $0{,}0673$. Parmi les rejets, la proportion de défauts est $2850/6730=285/673\simeq42{,}35\,\%$. Parmi les acceptations, elle est $150/93270=5/3109\simeq0{,}161\,\%$.
Deux contrôles peuvent réagir aux mêmes caractéristiques et être dépendants. Même avec les mêmes performances marginales, on ne peut pas multiplier sans connaître les probabilités conditionnelles pertinentes. Une hypothèse d'indépendance des résultats \emph{conditionnellement à l'état de la pièce} permettrait de calculer les chemins successifs ; elle doit être justifiée par le dispositif.

\section*{Problème 6 -- Un atelier de preuve et de données}
La somme des colonnes donne $339\,161$ femmes et $342\,812$ hommes, donc $681\,973$ personnes. Si $F_i$ est le nombre de femmes et $N_i>0$ la population du canton i, alors, avec $N=\sum_iN_i$,
\[
\sum_i\frac{N_i}{N}\frac{F_i}{N_i}
=\frac{\sum_iF_i}{N}.
\]
Les poids sont positifs et leur somme vaut un. Une moyenne simple accorderait autant de poids à Vianden qu'au canton de Luxembourg, malgré des effectifs très différents.
Le tableau attendu est Esch : $(97\,721;97\,370)$, autres : $(241\,440;245\,442)$. Les marges sont $(195\,091;486\,882)$ pour les lignes et $(339\,161;342\,812)$ pour les colonnes. Le module Python joint lit le CSV et construit ce tableau ; les tests vérifient ces sommes.
Si $u\le v$, écrire $tu+(1-t)v=u+(1-t)(v-u)$. Comme $0\le1-t\le1$, on obtient $u\le tu+(1-t)v\le v$. L'autre ordre s'obtient en échangeant u et v. Par le même argument, une moyenne pondérée positive se situe entre le minimum et le maximum de ses valeurs.

\end{document}
